% BT1201 -- Holonet lambda-lock theorem insert
% Generated from the page-36+ photonic-holonet consolidation pass.

\subsection{The \texorpdfstring{$\lambda$}{lambda}-Lock: drive, protection, helicity}
\label{sec:bt1201-lambda-lock}

The preceding self-correction paragraph contains one of the most useful
compression laws in the holonet: the same substrate deficit
\[
  \lambda=q-1
\]
appears in the drive, the topological protection, and the physical carrier.  The
safe statement is not that photon helicity generalises with the qudit dimension;
it does not.  Rather, the architecture is locked at the one substrate value for
which the general qudit pump law and the fixed massless-vector helicity count
coincide.

\begin{theorem}[The holonet $\lambda$-Lock]
\label{thm:bt1201-lambda-lock}
Let the finite substrate be the $q=3$ solution of the master equation
$q!=2q$, with $\lambda=q-1=2$.  Then three independently defined quantities
coincide:
\[
  \cos\theta_{\rm BC}=-\frac{q-1}{q}=-\frac23,
  \qquad
  |C|_{\max}=q-1=2,
  \qquad
  N_{\gamma,\perp}=2=\lambda.
\]
Here $\theta_{\rm BC}$ is the Boerdijk--Coxeter drive angle, $|C|_{\max}$ is the
extremal Chern protection of the spin-$(q-1)/2$ two-tone pump, and
$N_{\gamma,\perp}$ is the number of transverse helicities of a massless photon.
Thus
\[
  \boxed{
  q=3
  \quad\Longrightarrow\quad
  q-1=2
  =\text{drive deficit}
  =\text{Chern protection}
  =\text{photon helicity count}.
  }
\]
\end{theorem}

\begin{proof}[Proof and scope]
The BC drive angle is already fixed by the regular $q$-simplex relation
$\cos\theta_{\rm BC}=-(q-1)/q$, hence gives $-2/3$ at $q=3$.  The computed
spin-$S$ pump law gives extremal band Chern magnitude $2S$; for a dimension-$q$
carrier, $S=(q-1)/2$, hence $|C|_{\max}=q-1$.  The verifier checks this for
$q=2,3,4$, giving extremal magnitudes $1,2,3$.  Independently, a massless
spin-one gauge boson has exactly two physical transverse helicities.  The
coincidence with $q-1$ is therefore not a variable-helicity law for arbitrary
qudits; it is the compatibility condition that makes the $q=3$ substrate a
photonic substrate.
\end{proof}

This sharpens the architecture claim.  The holonet does not optimise over
qudit dimensions and then choose a photon.  It inherits $q=3$ from the substrate;
that inheritance forces $\lambda=2$; and $\lambda=2$ is exactly the point where
BC drive arithmetic, spin-one topological pumping, and massless-photon
polarisation all meet.  The qutrit is not a design preference.  It is the unique
place where the carrier and substrate agree.

\paragraph{Honesty boundary.}
The first two equalities have a natural $q$-family: the simplex drive deficit and
the spin-$(q-1)/2$ pump law.  The third equality is physical and specific:
ordinary photons have two transverse helicities because they are massless
spin-one gauge quanta.  The paper should therefore use the helicity equality as
a selection lock at $q=3$, not as a claim that a dimension-$q$ carrier possesses
$q-1$ helicities.
